\unit{Exposé 212 -- Quotient preschemes}
\sid{212}
\src{99}
\seminar{13th year, 1960/61, no. 212\\February 1961}
\paperhead{CONSTRUCTION TECHNIQUES AND EXISTENCE THEOREMS IN ALGEBRAIC GEOMETRY\\
III. QUOTIENT PRESCHEMES}%
  {by Alexander GROTHENDIECK}

\subsection*{Introduction}
\addcontentsline{toc}{subsection}{Introduction}

The problems treated in the present exposé differ from those considered in
the preceding two in that one seeks to represent certain covariant, rather
than contravariant, functors of variable schemes. The procedure of passing to
a quotient is nevertheless essential in many construction problems in
algebraic geometry, including those of Exposés I and II ([1], [2]). Thus the
question of the \emph{effectivity of descent data} on a $T$-prescheme $X$,
relative to a faithfully flat and quasi-compact morphism
$T\longrightarrow S$, is equivalent to the question of the existence of a
quotient of $X$ (satisfying the reasonable properties examined below) by the
flat equivalence relation in $X$ defined by the descent data; the questions
raised in [1], A, \S~2 c will no doubt be resolved at the same time as the
questions posed in paragraph 2 of the present exposé. Likewise, the
\emph{Picard scheme} of an $S$-scheme $X$ (for the definition, see [2], C,
\S~3) may be defined in various ways as a quotient of certain other schemes
(of positive divisors, or of immersions into a projective scheme) by flat
equivalence relations, while the definition and construction of these
auxiliary schemes is technically simpler: they are schemes of the form
$\underline{\operatorname{Hom}}_S(X,Y)$ and variants defined in [2], C,
\S~2, whose construction will be the subject of the next exposé (under
suitable projectivity hypotheses). Combining the results of the present
exposé and the next therefore yields the construction of Picard schemes under
suitable hypotheses.

The problem of passing to a quotient in the category of preschemes still
presents several unresolved questions. The most important is mentioned in
paragraph 8. It is at present the sole obstacle to the construction of
\emph{moduli schemes over the integers for curves of arbitrary genus,
polarized abelian varieties}, etc. This shows that its solution merits the
efforts of specialists in algebraic groups.

\src{100}
\subsection*{1. Equivalence relations and effective equivalence relations}
\addcontentsline{toc}{subsection}{1. Equivalence relations and effective equivalence relations}

Let $\mathcal C$ be a category and let $X$ be an object of $\mathcal C$. A
pair of morphisms
\[
  p_1,p_2:R\rightrightarrows X
\]
is called an \emph{equivalence pair} in $\mathcal C$, with \emph{target} $X$
and \emph{source} $R$, if for every object $T$ of $\mathcal C$ the two
corresponding maps
\[
  p_1(T),p_2(T):R(T)\rightrightarrows X(T)
\]
(where, for every object $Y$ of $\mathcal C$, we put
$Y(T)=\operatorname{Hom}(T,Y)$) define a map
\[
  R(T)\longrightarrow X(T)\times X(T)
\]
that induces a bijection from $R(T)$ onto the graph of an equivalence relation
on the set $X(T)$. There is an evident equivalence relation among the
equivalence pairs with target $X$; an equivalence class for this latter
relation is called a \emph{$\mathcal C$-equivalence relation on $X$}, or
simply an equivalence relation when no confusion can result.

If $X\times X$ exists, giving an equivalence relation on $X$ is equivalent to
giving a subobject $R$ of $X\times X$ such that, for every object $T$ of
$\mathcal C$, the subset of
\[
  (X\times X)(T)=X(T)\times X(T)
\]
corresponding to $R(T)$ is the graph of an equivalence relation on $X(T)$.
Writing $p_1$ and $p_2$ for the morphisms from $R$ to $X$ induced by the
projections $\operatorname{pr}_1$ and $\operatorname{pr}_2$, the preceding
condition says that $(p_1,p_2)$ is an equivalence pair. One may also express
diagrammatically in $\mathcal C$ (provided that $X\times X$ and the fiber
product $(R,p_2)\times_X(R,p_1)$ exist) the set-theoretic axioms of an
equivalence relation for the $R(T)$ in the $X(T)$, in accordance with the
general principle explained in [2], A, \S~1. We shall not need this.

Whenever there is a pair of morphisms $(p_1,p_2)$ with the same source $R$ and
the same target $X$, one may define the \emph{cokernel} of the pair as the
object $Y$ of $\mathcal C$ that represents the functor covariant in $Z$
\[
  \operatorname{Hom}_{p_1,p_2}(X,Z),
\]
the set of morphisms $u$ from $X$ to $Z$ such that $up_1=up_2$. If $Y$
exists, it is determined up to a unique isomorphism. We denote it by
$X/(p_1,p_2)$ or, by abuse of notation, by $X/R$, the latter notation being
used chiefly when
\src{101}$(p_1,p_2)$ is an equivalence pair: it is then customary to identify
in the notation the equivalence relation defined by the pair with $R$. Notice
that if $Y$ is regarded as a quotient of $X$, it in fact depends only on the
equivalence relation defined by the pair $(p_1,p_2)$.

Now start with a morphism
\[
  f:X\longrightarrow Y,
\]
which allows us to regard $X$ as an ``object over $Y$'', and suppose that the
fiber product
\[
  R(f)=X\times_Y X
\]
exists. Let $p_1$ and $p_2$ be its two projections. Then $(p_1,p_2)$ is an
equivalence pair, called the pair \emph{associated} with the morphism $f$. It
therefore defines an equivalence relation, called the relation associated with
$f$.

A pair of morphisms $(p_1,p_2)$ with target $X$ and common source $R$ is
called an \emph{effective equivalence pair} if
\begin{enumerate}
  \item[(i)] the cokernel $X/(p_1,p_2)=Y$ exists;
  \item[(ii)] the fiber product $X\times_Y X$ exists;
  \item[(iii)] the morphism $R\longrightarrow X\times_Y X$ with components $p_1$
  and $p_2$ is an isomorphism.
\end{enumerate}
Then $(p_1,p_2)$ is indeed an equivalence pair. The equivalence relation it
defines is also called an \emph{effective equivalence relation}.

A morphism $f:X\longrightarrow Y$ is called an \emph{effective epimorphism}
if
\begin{enumerate}
  \item[(i)] the fiber product $X\times_Y X=R$ exists;
  \item[(ii)] the quotient $X/(p_1,p_2)$ exists, where $p_1$ and $p_2$ are the two
  projections from $R$ to $X$;
  \item[(iii)] the morphism $X/(p_1,p_2)\longrightarrow Y$ induced by $f$ is an
  isomorphism.
\end{enumerate}
Then $f$ is indeed an epimorphism, and even a strict epimorphism (see [1], A,
\S~2.3); the converse holds if the fiber product $X\times_Y X$ exists. The
quotient object of $X$ defined by the epimorphism $f$ is also called an
\emph{effective quotient} of $X$.

The preceding definitions imply the following ``Galois correspondence'':

\src{102}
\result{PROPOSITION}{1.1} There is a one-to-one correspondence, preserving
the natural orders, between the set of effective equivalence relations $R$ on
$X$ and the set of effective quotients $Y$ of $X$: to $R$ corresponds the
effective quotient $X/R$, and to $Y$ corresponds the effective equivalence
relation defined by the canonical projection $X\longrightarrow Y$ (which is
defined by the fiber product $X\times_Y X$ equipped with its two projections).

In very good categories (sets, sheaves of sets, etc.), every quotient is
effective and every equivalence relation is effective. This is no longer so
in categories such as the category of preschemes over a fixed prescheme $S$,
even when $S$ is the spectrum of a field, and even when one restricts to
schemes finite over $S$. Questions of effectivity, and even (for preschemes
not finite over $S$) questions of the existence of quotients, most often prove
delicate.

\subsection*{2. Example: preschemes finite over $S$}
\addcontentsline{toc}{subsection}{2. Example: preschemes finite over S}

Let $\mathcal C$ be the category of preschemes finite over $S$, where $S$ is
assumed locally Noetherian. It is equivalent to the opposite of the category
of coherent sheaves of commutative algebras on $S$, or, when $S$ is affine
with ring $A$, to the opposite of the category of $A$-algebras finite over
$A$ (i.e. finitely generated as $A$-modules). It follows at once that finite
projective limits and finite inductive limits exist in $\mathcal C$. This is
well known (without any finiteness hypothesis\ldots) for the former. Thus the
fiber product of preschemes $X,Y$ over $S$ corresponds to the tensor product
$B\otimes_A C$ of the corresponding algebras; the kernel of two morphisms
$X\rightrightarrows Y$, defined by two homomorphisms of $A$-algebras
$u,v:C\rightrightarrows B$, corresponds to the quotient of $B$ by the ideal
generated by the elements $u(c)-v(c)$, etc. For finite inductive limits it is
enough to consider, on the one hand, finite sums, which correspond to the
ordinary product of $A$-algebras, and, on the other hand, cokernels of pairs of
morphisms $X\rightrightarrows Y$, which correspond (as one sees immediately)
to the subring of $C$ consisting of the elements on which the homomorphisms
$u,v:C\rightrightarrows B$ coincide (this subring is finite over $A$ by the
Noetherian hypothesis). One may moreover show, using the Noetherian
hypothesis, that the finite inductive limits, and in particular the quotients,
so constructed in the category $\mathcal C$ of preschemes finite over $S$ are
in fact quotients in the category of \emph{all} preschemes.

\src{103}
As we said in [1], there are non-effective epimorphisms in $\mathcal C$ (or
non-strict epimorphisms\footnote{The printed text omits ``non'' before
``strict''. In the category under consideration fiber products exist, and the
preceding paragraph establishes precisely the equivalence between effectivity
and strict epimorphy.}; this is the same thing since fiber products exist).
\emph{I do not know whether equivalence relations are always effective} when
no flatness hypothesis is imposed. In this direction I have obtained only very
partial positive results, indispensable for the proof of the fundamental
theorem in the formal theory of modules (see [2], B, th. 1). Let us point out
that in the problem posed it is easy to reduce to the case where $S$ is the
spectrum of a local Artinian ring $A$ with algebraically closed residue field.
But even if $A$ is a field, the answer is not known.

One may also consider a prescheme $X$ over $S$ that is no longer assumed
finite over $S$, but with an equivalence relation $R$ on $X$ such that
$p_1:R\longrightarrow X$ is a finite morphism. Then $R$ is called a finite
equivalence relation. Suppose for simplicity that $S$ and $X$ are affine (and
hence $R$ is affine, so that the situation reduces to one of pure commutative
algebra). \emph{Even in this case it is not known whether a quotient $X/R=Y$
exists and whether the canonical morphism $X\longrightarrow Y$ is finite}.
(The simplest case is that in which $S$ is the spectrum of a field $k$ and
$X$ is the spectrum of $k[t]$, i.e. the affine line.) Of course, if the two
preceding problems have affirmative solutions, one may conclude in the present
situation that $R$ is effective. Notice that the problem of the
\emph{existence} of a quotient $Y$ and the finiteness of
$f:X\longrightarrow Y$ arises in exactly the same terms if, instead of the
graph of an equivalence relation on $X$, one has only an equivalence pregraph
on $X$ in the sense of paragraph 4.

The question of passing to the quotient by a more or less arbitrary finite
equivalence relation arises in constructing preschemes by ``gluing'' given
preschemes $X_i$ along certain closed subpreschemes; the gluing law is
expressed precisely by a finite equivalence relation on the sum prescheme $X$
of the $X_i$. One should also expect the solution of the problems posed here,
and of various variants, to be a preliminary condition for developing a
general technique of nonprojective constructions in the direction initiated
in [2].

The only general positive fact known to the author is the following:

\result{PROPOSITION}{2.1} Let $S$ be a locally Noetherian prescheme, let $s$
be a point of $S$, and let $\Omega$ be an algebraically closed extension of
$k(s)$. Consider the
\src{104}corresponding ``fiber functor'' $F$, which assigns to every $S$-scheme
$X$ finite over $S$ the set of $\Omega$-valued points of $X_s$. This functor
(trivially left exact) is \emph{right exact}, i.e. it commutes with finite
inductive limits, and in particular with cokernels of pairs of morphisms.

Using this result for all the ``geometric points'' of $S$, one deduces that
the ``quotient'' category $\mathcal C'$ of $\mathcal C$, obtained by reasoning
``modulo surjective radicial morphisms'' (i.e. by formally adjoining inverses
to the morphisms in question), is a ``geometric'' category, i.e. it satisfies
the same properties ``of finite nature'' as the category of sets. In
particular, every equivalence relation in it is effective. This implies that
if $R$ is an equivalence relation on an $X$ finite over $S$, then the canonical
morphism
\[
  R\longrightarrow X\times_Y X\qquad (\text{where }Y=X/R)
\]
is \emph{radicial and surjective} (in fact a surjective closed immersion,
since it is a monomorphism).

\subsection*{3. The case of an operator group}
\addcontentsline{toc}{subsection}{3. The case of an operator group}

Assume again that $\mathcal C$ is an arbitrary category. Let $G,X$ be objects
of $\mathcal C$, and suppose that $G$ is a $\mathcal C$-operator group on the
object $X$ of $\mathcal C$. This means (see [2], A, \S~1) that for every
object $T$ of $\mathcal C$ a group structure has been given on $G(T)$ and a
$G(T)$-set structure on $X(T)$, in such a way that as $T$ varies the
structures in question ``vary functorially'' with $T$. If the products
$G\times G$ and $G\times X$ exist in $\mathcal C$, such a structure may also
be defined as a pair of morphisms
\[
  G\times G\longrightarrow G,
  \qquad \pi:G\times X\longrightarrow X,
\]
subject to the condition that for every object $T$ of $\mathcal C$, the
corresponding laws of composition on the sets $G(T)$ and $X(T)$ make $G(T)$ a
group acting on $X(T)$. Translating this axiom into the commutativity of
certain diagrams in $\mathcal C$ is easy but tedious, and in fact completely
useless in every case known to me.

Suppose that $G\times X$ exists, and consider the two morphisms
\[
  p_1,p_2:G\times X\rightrightarrows X
\]
with
\src{105}
\[
  p_1=\operatorname{pr}_2,\qquad p_2=\pi.
\]
One sees immediately that $(p_1,p_2)$ is an equivalence pair if and only if,
for every object $T$ of $\mathcal C$, the map
\[
  G(T)\times X(T)\simeq(G\times X)(T)
  \longrightarrow X(T)\times X(T),
\]
defined by this pair is injective, i.e. if the group $G(T)$ acts
\emph{freely} on the set $X(T)$, that is, if $g\in G(T)$, $x\in X(T)$, and
$g\cdot x=x$ imply that $g$ is the identity element of $G(T)$. One then says
that $G$ \emph{acts freely on $X$} (or that $X$ is a
\emph{$\mathcal C$-principal space under $G$}). The equivalence relation
associated with the pair $(p_1,p_2)$ is then called the \emph{equivalence
relation defined by the group $G$ acting freely on $X$}. If $X\times X$ also
exists and one considers the morphism
\[
  p:G\times X\longrightarrow X\times X
\]
defined by the pair $(p_1,p_2)$, the condition that $G$ act freely means that
$p$ is a \emph{monomorphism}.

Of course, even when $G$ does not act freely on $X$, one wants criteria for
the existence of a quotient of $X$ by $G$, i.e. of the cokernel of the
preceding pair $(p_1,p_2)$.

The cokernel in question will often be denoted by $X/G$, or preferably by
$G\backslash X$ if $G$ acts on the left (the former notation being reserved
for the case in which $G$ acts on the right). Notice that even when the
``image'' of $G\times X$ under $p$ exists (this image being defined, for
example, as the smallest subobject of $X\times X$ through which $p$ factors),
say $R$, it is most often not an equivalence relation on $X$. If one then
tries to pass directly to the quotient by $R$ (or more precisely by the pair
of morphisms from $R$ to $X$ induced by the two projections
$\operatorname{pr}_i$), one irretrievably loses the special features of the
original pair $(p_1,p_2)$. It is therefore important to find a generalization
of the notion of equivalence relation that applies directly to the pair
defined by a $\mathcal C$-operator group.

\subsection*{4. Pre-equivalence relations}
\addcontentsline{toc}{subsection}{4. Pre-equivalence relations}

Recall that a category in which every morphism is an isomorphism is called a
\emph{groupoid}. On the other hand, a category must be defined as consisting
of two underlying sets $(X,R)$, the set of \emph{objects} and the set of
\emph{arrows}, equipped with the following structures:
\src{106}
\begin{enumerate}
\item[(i)] A pair of maps
\[
  p_1,p_2:R\rightrightarrows X
\]
called the \emph{source map} and the \emph{target map},

\item[(ii)] A map
\[
  \pi:(R,p_2)\times_X(R,p_1)\longrightarrow R
\]
called the \emph{composition map}.
\end{enumerate}
These data must satisfy familiar axioms, which we shall not repeat here and
which could be expressed by the commutativity of certain diagrams, together
with the existence of a map $D:X\to R$ (necessarily unique) making two other
diagrams commute. The map $D$ sends an object to its corresponding identity
arrow and satisfies
\[
  p_1D=p_2D=\operatorname{id}_X.
\]
To say that the category is a groupoid then means that there exists a map
\[
  s:R\longrightarrow R
\]
(necessarily unique), called the \emph{inversion} of $R$, which transforms
each arrow into an inverse arrow. This could be expressed by the commutativity
of four further diagrams formed from $s$, $\Delta$, and the preceding data;
the first two are
\[
  p_1s=p_2,\qquad p_2s=p_1.
\]

With these points recalled, the general definitions in ([2], A, \S~1) show,
in particular, what is meant by a structure of a $\mathcal C$-\emph{category},
respectively a $\mathcal C$-\emph{groupoid}, on a pair of objects $X,R$ of an
arbitrary category $\mathcal C$: by definition, it is the datum, for every
object $T$ of $\mathcal C$, of a category structure, respectively a groupoid
structure, in the set-theoretic sense whose set of objects is $X(T)$ and whose
set of arrows is $R(T)$, these structures ``varying functorially'' with the
variable $T$. This therefore gives two morphisms
\[
  p_1,p_2:R\rightrightarrows X
\]
called the \emph{source morphism} and the \emph{target morphism}, and, when
the required fiber product exists, a morphism
\src{107}
\[
  \pi:(R,p_2)\times_X(R,p_1)\longrightarrow R
\]
called the \emph{composition morphism}. These three data then suffice to
determine the category structure on $X,R$, the axiom imposed on them being as
follows: for every $T$, the three corresponding data for $X(T),R(T)$ define a
category structure (respectively a groupoid structure) on this pair of sets.
When appropriate, this may be expressed by the commutativity of certain
diagrams involving a uniquely determined morphism
\[
  D:X\longrightarrow R
\]
and, in the case of groupoids, a uniquely determined morphism
\[
  s:R\longrightarrow R;
\]
the diagrams are written out just as in the ``set-theoretic'' case. This
laborious interpretation of the axioms is fortunately useless in practice.
The sole theoretical interest of being able to express the data and axioms in
terms of morphisms and equalities of morphisms between certain fiber products
is the following: if $F:\mathcal C\to\mathcal C'$ is a left exact functor
(i.e. one commuting with finite products and fiber products), then it
transforms a $\mathcal C$-category (respectively a $\mathcal C$-groupoid) into
a $\mathcal C'$-category (respectively a $\mathcal C'$-groupoid), provided
finite products and fiber products exist in $\mathcal C$.

In practice, it is important to know how to interpret the morphisms
$p_1,p_2,\pi,D,s$ as \emph{simplicial operations} in a suitable semisimplicial
or simplicial object of $\mathcal C$, at least when fiber products exist in
$\mathcal C$. To fix terminology, let us introduce the \emph{category
$\mathcal S$ of standard simplices} as the category whose objects are the
finite sets of the form
\[
  \Delta_n=[0,n],\qquad\text{where }n\geq 0
\]
(the interval of integers from $0$ to $n$), and whose morphisms or arrows are
\emph{arbitrary maps} between these finite sets. Notice that the category
$\mathcal S$ is equivalent to the category of \emph{nonempty} finite sets,
with maps of finite sets as morphisms. In $\mathcal S$, the coproduct of a
finite \emph{nonempty} family of objects clearly exists, as does the pushout
of two objects under a third (the operation dual to the fiber product). Let
$\mathcal S'$ denote the subcategory of $\mathcal S$ having the same objects
but whose morphisms are the \emph{increasing maps} between the $\Delta_n$.
This category is equivalent to the category of nonempty totally ordered
finite sets. In this category the
\src{108}
coproduct of two objects never exists, and the pushout of two objects $A,B$
under a third $C$ does not exist in general (take, for example,
$C=\Delta_0$, $A=B=\Delta_1$, with the two structure maps $u:C\to A$ and
$v:C\to B$ equal). It does exist in certain cases, however; for example,
\[
  C=\Delta_0,\qquad A=\Delta_m,\qquad B=\Delta_n,\qquad
  u(0)=m,\qquad v(0)=0,
\]
and in this case
\[
  \Delta_m\amalg_{\Delta_0}\Delta_n=\Delta_{m+n}.
\]

A \emph{simplicial object} (respectively a \emph{semisimplicial object}) in a
category $\mathcal C$ is, by definition, a contravariant functor $K$ from
$\mathcal S$ (respectively from $\mathcal S'$) to $\mathcal C$. Thus a
simplicial object defines a semisimplicial object by restriction, but the
first notion differs from the second essentially through the presence of
\emph{symmetry operations} on the $K_n=K(\Delta_n)$, corresponding to the
images under the functor $K$ of the elements of the symmetric group on $n+1$
letters (regarded as the automorphism group of $\Delta_n$ in $\mathcal S$).

For every $n$, let $\Delta'_n$ (respectively $\Delta''_n$) be the finite
category whose object set is $\Delta_n$, defined by the chaotic order relation
(respectively the natural total order) on $\Delta_n$ (the set of arrows being
the graph of that order relation). It is clear that $\Delta'_n$ (respectively
$\Delta''_n$) depends functorially on the object $\Delta_n$ of $\mathcal S$
(respectively $\mathcal S'$). If $Z$ is a category, then
$\operatorname{Hom}(\Delta'_n,Z)$ (respectively
$\operatorname{Hom}(\Delta''_n,Z)$), as $\Delta_n$ varies, is a functor from
$\mathcal S$ (respectively $\mathcal S'$) to the category of sets, i.e. a
simplicial set (respectively a semisimplicial set), said to be
\emph{associated with the category $Z$}; denote them by $Z'$ and $Z''$. There
is moreover an evident natural homomorphism from the semisimplicial set
associated with $Z'$ to $Z''$, and the latter is an isomorphism if and only if
$Z$ is a groupoid. With this understood:

\result{PROPOSITION}{4.1} The functor $Z\longmapsto Z''$ from the category of
categories to the category of semisimplicial sets is fully faithful, and
defines an equivalence between the category of \emph{categories} and the
category of semisimplicial sets, i.e. contravariant functors $K$ from
$\mathcal S'$ to $(\mathrm{Set})$, which transform pushouts $A\amalg_C B$ of
the type specified above into fiber products of sets. Likewise, the functor
$Z\longmapsto Z'$ from the category of groupoids to the category of
simplicial sets is fully faithful, and defines an equivalence between the
category of \emph{groupoids} and the category of
\src{109}
simplicial sets, i.e. contravariant functors $K$ from $\mathcal S$ to
$(\mathrm{Set})$, which transform pushouts into fiber products.

Thus categories may be regarded as particular semisimplicial sets, and
groupoids as particular simplicial sets, provided of course that one reasons
``up to isomorphism,'' as is required when interpreting certain structures in
terms of others. The usual procedure of reduction to the set-theoretic case
then implies:

\result{COROLLARY}{4.2} The preceding statement remains valid when categories,
groupoids, and simplicial sets are replaced by $\mathcal C$-categories,
$\mathcal C$-groupoids, and simplicial objects in $\mathcal C$, provided that
fiber products exist in $\mathcal C$.

The semisimplicial object $K$ in $\mathcal C$ associated with a category
$(X,R,\ldots)$ in $\mathcal C$ may be made explicit by interpreting the
component $K_n=K(\Delta_n)$ of $K$ as the $n$-fold fiber power of
$(R,p_1)$ over $X$\footnote{The printed text says
``$(n+1)$-fold fiber power'' and, in the following formula, $K_0=R$. The
object component of the nerve is $K_0=X$ and its arrow component is $K_1=R$;
the two index corrections are therefore linked.}, or better, by the
recurrence formula
\[
  K_0=X,\qquad
  K_n=(K_{n-1},p_{n-1}^{(n-1)})\times_X(R,p_1),
\]
where the $p_i^{(n-1)}$ ($0\leq i\leq n-1$) are the natural projections from
$K_{n-1}$ to $X$ (which are themselves also defined recursively). In this
way, $p_1,p_2,\pi,D,s$ are interpreted as the simplicial operations
corresponding respectively to the following morphisms in $\mathcal S$: face
$0$ of $\Delta_1$, face $1$ of $\Delta_1$, face $(0,2)$ of $\Delta_2$, the
degeneracy $\Delta_1\to\Delta_0$, and the symmetry of $\Delta_1$. All other
semisimplicial (respectively simplicial) operations are obtained formally from
the preceding four (respectively five) by composition and fiber products.

A \emph{pre-equivalence relation} on an object $X$ of a category is now
defined to be the datum of a groupoid whose object of objects (so to speak) is
$X$. Such a datum therefore includes, among other things, an object $R$ and
two morphisms
\[
  p_1,p_2:R\rightrightarrows X.
\]
But notice that these data alone do not determine the structure in question,
in contrast to what happens for equivalence pairs. In this exposé, we are
interested in this notion in order to obtain criteria for the possibility of
passing to the quotient, i.e. for forming the cokernel of the pair
$(p_1,p_2)$. Thus the statement of the problem does not invoke the additional
data inherent in a groupoid. In proving the results that follow, we shall
\src{110}
nevertheless use these additional data, and in particular the simplicial
operations (including symmetry operations) through dimension $3$ (involving
the fourfold fiber power of $R$ over $X$).

An equivalence relation on an object $X$ of $\mathcal C$ defines a
pre-equivalence relation: it suffices to see this in the set-theoretic case,
where one associates with an equivalence relation on a set $X$ the groupoid
whose set of objects is $X$ and whose set of arrows is the graph of the
equivalence relation.

A $\mathcal C$-monoid $G$ acting on an object $X$ of $\mathcal C$ defines a
$\mathcal C$-category whose underlying objects are $R=G\times X$ and $X$
(provided that $G\times X$ exists), and this is a $\mathcal C$-groupoid if and
only if $G$ is a group. Again, it suffices to verify this in the set-theoretic
case. The composition of the arrows $(g,a)$ and $(g',g\cdot a)$ is then
defined by
\[
  (g',g\cdot a)(g,a)=(g'g,a),
\]
i.e. if $a,b\in X$, then $\operatorname{Hom}(a,b)$ is by definition the
transporter from $a$ to $b$, and the morphisms compose by means of composition
of the elements of $G$.

\medskip
\noindent\textbf{REMARK. ---} The logical difficulties raised by a statement
such as (4.1) can be avoided by understanding that all the objects under
consideration lie in a fixed ``Universe'' (which is itself a set).

\fsection{5}{Quotient by a finite flat equivalence relation}

\result{THEOREM}{5.1} Let $X=\operatorname{Spec}(B)$ be an affine scheme and
$R$ a pre-equivalence relation on $X$ whose component $R_1$ is affine, say
$R_1=\operatorname{Spec}(C)$. Suppose that the first projection
$p_1:R_1\to X$ is finite and locally free; in other words, suppose that the
corresponding ring homomorphism $p_1^*:B\to C$ makes $C$ a finitely generated
projective $B$-module. Let $A$ be the subring of $B$ that is the kernel of the
pair of homomorphisms
\[
  p_1^*,p_2^*:B\rightrightarrows C
\]
(the set of elements $b$ such that $p_1^*(b)=p_2^*(b)$). Let
$Y=\operatorname{Spec}(A)$, and let $f:X\to Y$ be the morphism defined by the
inclusion of $A$ into $B$. Under these conditions:
\begin{enumerate}
\item[(i)] $B$ is integral over $A$, i.e. $f$ is an integral morphism.

\item[(ii)] The morphism $f$ is surjective, and its fibers are the equivalence
classes
\src{111}
$p_2(p_1^{-1}(x))$ in the set $X$ modulo $R$; moreover, the topology of $Y$ is
the quotient of that of $X$.

\item[(iii)] $Y$ is the quotient of $X$ by $R$ in the category of preschemes.

\item[(iv)] If $R$ comes from an \emph{equivalence relation}, then the
morphism $f:X\to Y$ is finite locally free (i.e. $B$ is a finitely generated
projective $A$-module), and the equivalence relation is effective, i.e.
$R_1\to X\times_Y X$ is an isomorphism.
\end{enumerate}

This theorem generalizes the familiar theorem concerning a finite group $G$
acting by automorphisms on the ring $B$ and the invariant ring $A$, and its
proof is analogous to the familiar proof. Statement (iii) can be sharpened as
follows:

\result{COROLLARY}{5.2} The canonical morphism
$R_1\to X\times_Y X$ is \emph{surjective}.

Let $R$ still be a ``finite and locally free'' pre-equivalence relation on
$X$, but now let $X$ be an arbitrary prescheme. Suppose that one can find a
prescheme $Y$ and a morphism $f:X\to Y$ such that $fp_1=fp_2$ and, moreover,
such that the sequence of homomorphisms of sheaves of rings on $Y$
\[
  \mathcal O_Y\longrightarrow f_*(\mathcal O_X)
  \rightrightarrows g_*(\mathcal O_R)
\]
is exact (where $g=fp_i$). The theorem then gives conclusions analogous to
(i)--(iv) above; in particular, by (iii), $Y$ is the quotient of $X$ by $R$
and is therefore determined up to unique isomorphism. Under these conditions
the pre-equivalence relation $R$ on $X$ is said to be \emph{admissible}. With
this definition:

\result{THEOREM}{5.3} Let $X$ be a prescheme and $R$ a pre-equivalence relation
on $X$ such that $p_1:R_1\to X$ is finite and locally free. Then $R$ is
admissible if and only if every set-theoretic equivalence class
$p_2(p_1^{-1}(x))$ in $X$ modulo $R$ is contained in an affine open subset
(a condition that is always satisfied if every finite subset of $X$ is
contained in an affine open subset, for example if $X$ is quasi-projective
over an affine scheme).

Indeed, one shows without difficulty that every equivalence class modulo $R$
in $X$ is then contained in an affine open subset \emph{stable under $R$}, and
one constructs the quotient $Y$ by gluing the pieces obtained from Theorem
5.1.
\src{112}

\result{COROLLARY}{5.4} Suppose that this condition is satisfied and that, in
addition, $R$ comes from an equivalence relation. Then that relation is
effective, i.e. $R_1\to X\times_Y X$ is an isomorphism, and $f:X\to Y$ is
finite and locally free.

One immediately concludes by ``descent'':

\result{COROLLARY}{5.5} Under the hypotheses of (5.4), $X$ has rank $n$
everywhere over $Y$ if and only if $(R_1,p_1)$ has rank $n$ everywhere over
$X$. If $X,R_1$ are $Z$-preschemes and $p_1,p_2$ are $Z$-morphisms, so that
$Y$ is a $Z$-prescheme, then $X$ is flat over $Z$ if and only if $Y$ is flat
over $Z$.

In summary:

\result{SCHOLIUM}{} The datum of a finite, locally free, and surjective
morphism $f:X\to Y$ of preschemes is equivalent to the datum of a prescheme
$X$ equipped with an equivalence relation $R$ such that $p_1:R\to X$ is
finite and locally free and every class $p_2(p_1^{-1}(x))$ is contained in an
affine open subset.

\medskip
\noindent\textbf{REMARKS 5.6.}

\smallskip
\noindent $1^\circ$ No noetherian hypothesis was needed.

\smallskip
\noindent $2^\circ$ This notion of passage to the quotient includes, as a
special case, CARTIER's ``inseparable descent,'' which amounts to determining
the finite locally free morphisms $f:X\to Y$ such that $f_*(\mathcal O_X)$
admits a $p$-basis relative to $\mathcal O_Y$ (where $X$ is a given prescheme
whose sheaf $\mathcal O_X$ is annihilated by the prime number $p>0$). Notice
that the result is readily expressed without any regularity hypothesis on the
local rings and without assuming that $X$ is an algebraic scheme over a field.
The JACOBSON--BOURBAKI theory is obtained by taking $X$ to be the spectrum of
a field of characteristic $p$.

\smallskip
\noindent $3^\circ$ GABRIEL had previously obtained a special case of Theorem
(5.3) in the theory of passage to the quotient for finite commutative groups
over a field $k$. [Compare (7.4).]

\fsection{6}{Quotient by a proper flat equivalence relation}

\result{THEOREM}{6.1} Let $S$ be a locally noetherian prescheme, $X$ a
quasi-projective $S$-scheme, and $R$ a pre-equivalence relation on $X$ such
that:
\src{113}
\begin{enumerate}
\item[(a)] $p_1:R_1\to X$ is proper and flat;
\item[(b)] $R_1\to X\times_S X$ is finite, or equivalently, by (a), has
finite fibers (a condition automatically satisfied if $R$ comes from an
equivalence relation).
\end{enumerate}
Under these conditions:
\begin{enumerate}
\item[(i)] $X/R=Y$ exists and, if $S$ is noetherian and $R$ comes from an
equivalence relation, is quasi-projective over $S$.\footnote{The printed text
does not contain the second hypothesis. The linked erratum to Exposé 212 says
that quasi-projectivity then appears to have been proved only when $R$ comes
from an equivalence relation.}

\item[(ii)] The canonical morphism $f:X\to Y$ is surjective, proper, and open;
its fibers are the equivalence classes $p_2(p_1^{-1}(x))$ in $X$ modulo $R$.
Thus $Y$ is identified with the quotient topological space of $X$ by the
set-theoretic equivalence relation defined by $R$. Finally,
$R_1\to X\times_Y X$ is surjective.

\item[(iii)] If $R$ comes from an equivalence relation, then that relation is
effective, i.e. $R_1\to X\times_Y X$ is an isomorphism, and moreover
$f:X\to Y$ is flat (hence faithfully flat).
\end{enumerate}

For the proof, one reduces to (5.1) by means of suitable quasi-sections of $X$
for $R$, the proof being analogous to the construction of quotient algebraic
groups in the Chevalley Seminar.

In summary:

\result{SCHOLIUM}{} Let $X$ be quasi-projective over a locally noetherian
$S$. The datum of a proper, faithfully flat, and surjective morphism
$f:X\to Y$ from $X$ to an $S$-prescheme $Y$ is equivalent to the datum of an
equivalence relation $R$ on $X$ such that $p_1:R\to X$ is proper and flat.

The same method gives the following result:

\result{THEOREM}{6.2} Let $S$ be a noetherian prescheme, $X$ a prescheme of
finite type over $S$, and $R$ a pre-equivalence relation on the $S$-prescheme
$X$. Suppose that:
\begin{enumerate}
\item[(a)] $p_1:R_1\to X$ is flat and of finite type;
\item[(b)] the morphism $R_1\to X\times_S X$ is quasi-finite (i.e. has finite
fibers).
\end{enumerate}
Then there exists a dense open subset $U$ of $X$ that is \emph{saturated} for
$R$, such that:
\begin{enumerate}
\item[(i)] If $R_U$ is the pre-equivalence relation induced by $R$ on $U$,
then $U/R_U$ exists and is of finite type over $S$.

\item[(ii)] The canonical morphism $U\to U/R_U$ is surjective and open; its
fibers are the set-theoretic equivalence classes for $R_U$ (thus $U/R_U$ is
the quotient topological space of $U$ by the set-theoretic equivalence
relation defined by $R_U$). Finally, the morphism
\[
  (R_U)_1\longrightarrow U\times_{U/R_U}U
\]
is surjective.
\src{114}

\item[(iii)] If $R$ comes from an equivalence relation, one may arrange for
$U\to U/R_U$ to be faithfully flat and for $R_U$ to be effective.
\end{enumerate}

This is a result of an essentially ``birational'' nature.

\medskip
\noindent\textbf{REMARKS 6.3.}

\smallskip
\noindent $1^\circ$ I do not know whether hypothesis (b) is unnecessary in
(6.1) and (6.2). In taking quotients by groups, it forces us in practice to
restrict ourselves to the case in which all stabilizers are finite groups.

\smallskip
\noindent $2^\circ$ One may ask whether there are results analogous to (6.1)
and (6.2) without a flatness hypothesis. I have no counterexample in this
direction. On the other hand, even retaining the flatness hypothesis and
restricting to equivalence relations for which $p_1:R\to X$ is flat and
quasi-finite (but not finite), with $X$ affine, it can happen that $R$ is not
effective: take the equivalence relations induced on affine open subsets
covering Nagata's variety (on which a group of order two acts in a
``non-admissible'' manner).

\fsection{7}{Applications}

As we said in the introduction, the most important application of (6.1) is
the construction of Picard schemes, as well as solutions to various other
moduli problems, to which we shall return later.

One obtains a simple proof of the following result of SHIMURA:

\result{PROPOSITION}{7.1} Let $A$ be an abelian scheme defined over a discrete
valuation ring $V$ with fraction field $K$. Then every abelian scheme $B'$ over
$K$ that is isogenous to a quotient of $A\otimes_V K$ ``has good reduction
over $V$,'' i.e. is isomorphic to some $B\otimes_V K$,\footnote{The printed
text has $B\otimes_V A$. Since $B$ is a $V$-scheme and the assertion concerns
its generic fiber, the base factor must be $K$.} where $B$ is an abelian
scheme over $V$ (essentially unique, as we recall).

We may suppose that $B'$ is the quotient of $A_K$ by a subgroup scheme $C'$.
(N.B.: $C'$ will not in general be ``reduced,'' i.e. its local rings will have
nilpotent elements.) Consider the closed subscheme $C$ of $A$ that is the
``closure'' of $C'$, i.e. the smallest closed subscheme of $A$ such that $C_K$
contains $C'$. We have $C_K=C'$, and the fact that $V$ is a discrete valuation
ring readily implies that $C$ is a subgroup scheme of $A$ over $V$. Since $A$
is proper over $\operatorname{Spec}(V)=S$, so is $C$; moreover, $A$ is even
projective over $S$. We may therefore apply (6.1) to construct $A/C=B$,
\src{115}
which is the desired $B$.

Finally, essentially familiar arguments allow one to deduce from (6.2) the
following result:

\result{THEOREM}{7.2} Let $S$ be the spectrum of an artinian ring, let $F$ and
$G$ be group schemes of finite type over $S$, and let $u:F\to G$ be a
homomorphism of group schemes over $S$. Suppose that:
\begin{enumerate}
\item[(i)] $F$ is flat over $S$;
\item[(ii)] the kernel of $u$ is finite.
\end{enumerate}
Under these conditions, the quotient scheme $G/F$ exists; the canonical
morphism $G\to G/F$ is surjective and open, and its fibers are the
set-theoretic equivalence classes defined by the right actions of $F$ on $G$.
Finally, if $u$ is a monomorphism, then $G\to G/F$ is flat and the morphism
\[
  G\times F\longrightarrow G\times_{G/F}G
\]
is an isomorphism. In other words, $G$ is a principal homogeneous space over
$G/F$ with structure group (acting on the right) $F$, or rather
$F\times_S(G/F)$ regarded as a group scheme over $G/F$ (see [1], B, \S~6).

\result{COROLLARY}{7.3} Under these conditions, $G$ is flat over $S$ if and
only if $G/F$ is flat over $S$. If this condition holds, formation of the
quotient by $F$ commutes with every base extension of $S$; and if $F$ is a
normal subgroup of $G$, then $G/F$ can be endowed with the structure of the
\emph{quotient group} of $G$ by $F$.

The situation is particularly simple when $S$ is the spectrum of a field, for
then every $S$-prescheme is automatically flat over $S$. One obtains:

\result{COROLLARY}{7.4} Let $F,G$ be group schemes of finite type over a field
$k$, and let $u:F\to G$ be a homomorphism of $k$-groups. Then $u$ factors as
$F\to F'\to G$, where $F\to F'$ is the quotient homomorphism by the closed
subgroup $\operatorname{Ker}u$ of $F$, and $F'\to G$ is a group homomorphism
that is a closed immersion. The quotient $G/F=G/F'$ exists. The usual
formalism (of the kind provided by Noether's theorems) is valid among
algebraic groups over $k$.

This result makes it possible to treat uniformly both passage to the quotient
among algebraic groups in the classical sense (i.e. those which are
irreducible over $k$ and smooth over $k$) and passage to the quotient by
``infinitesimal'' subgroups
\src{116}
considered by CARTIER. It is advantageous to regard the ``hyperalgebras''
introduced by CARTIER, following DIEUDONNÉ's work on formal groups, as groups
in the category of formal schemes over $k$, and, when applicable (if they
correspond to hyperalgebras of finite rank over $k$), as finite algebraic
groups over $k$.

\fsection{8}{A conjecture}

It answers the need to form quotients by the projective group acting on
certain subschemes of ``Hilbert schemes'' (the latter replacing Chow varieties
in scheme theory).

Let $S$ be a prescheme and $n$ an integer. To every prescheme $S'$ over $S$,
associate the group $\operatorname{Gl}(n,\Gamma(S',\mathcal O_{S'}))$ of
invertible $n$-by-$n$ matrices with entries in the ring of sections of
$\mathcal O_{S'}$. This gives a contravariant functor in $S'$, which is readily
shown to be representable; it therefore corresponds to a group scheme over
$S$, in fact affine over $S$, denoted $\operatorname{Gl}(n)_S$. Its formation
is compatible with base change, so in fact everything comes from a group
scheme over $\mathbf Z$, denoted $\operatorname{Gl}(n)$. The group
$\operatorname{Gl}(1)$, called the multiplicative group and often denoted
$\mathbf G_m$, corresponds to the functor
$S\longmapsto\Gamma(S,\mathcal O_S)^*$, the group of ``units'' over $S$. There
is an evident homomorphism $\operatorname{Gl}(1)\to\operatorname{Gl}(n)$, and
one readily constructs the quotient group, denoted $\operatorname{GP}(n-1)$
and called the \emph{projective group of degree $n-1$ over $\mathbf Z$}. It
represents the functor associating to $S$ the group of sections of the sheaf
$\widetilde{\operatorname{Gl}(n)}_S/
\widetilde{\operatorname{Gl}(1)}_S$, where
$\widetilde{\operatorname{Gl}(n)}_S$ denotes the sheaf of germs of sections of
$\operatorname{Gl}(n)_S$ over $S$. (Care must be taken: sections of
$\operatorname{GP}(n-1)_S$ over $S$ do not in general come from sections of
$\operatorname{Gl}(n)_S$ over $S$!) Notice that one can prove that
$\operatorname{GP}(n-1)$ also represents the functor
$S\longmapsto\operatorname{Aut}_S(\mathbf P_S^{n-1})$, where
$\mathbf P_S^{n-1}$ is standard projective space of relative dimension $n-1$
over $S$, at least for noetherian $S$. This is how it enters moduli theory.

Let $S$ be a noetherian scheme, which may if desired be assumed affine, and
let $X$ be a quasi-projective $S$-prescheme equipped with an invertible sheaf
$\mathcal L$ that is very ample relative to $S$. Suppose that the group
$G=\operatorname{GP}(n)_S$ acts on $X$ and at the same time on $\mathcal L$
(compatibly with its action on $X$), and that it acts \emph{freely} on
$X$.\footnote{The printed text says ``on $S$.'' The action under consideration
and the quotient condition concern $X$.}
\src{117}

\result{CONJECTURE}{8.1}\footnote{The linked erratum to Exposé 212 states
that this conjecture is definitively false. It also reports that the positive
fact in Remark 8.1 was subsequently proved and refers to Mumford's theorem and
later refinements.} Under the preceding conditions:

\smallskip
\noindent $1^\circ$ \emph{The equivalence relation defined by $G$ is
effective, the quotient $Y=X/G$ is of finite type over $S$, and the canonical
morphism $f:X\to Y$ is flat and surjective} (thus $X$ becomes a principal
homogeneous bundle over $Y$ with group
$G\times_S Y=\operatorname{GP}(n)_Y$).

\smallskip
\noindent $2^\circ$ Let $\mathcal L'$ be the invertible sheaf on $Y$ deduced
from $\mathcal L$ by ``faithfully flat descent'' along $f$ (see [1], B,
\S~th. 1). Then $\mathcal L'$ is \emph{pre-ample on $Y$ relative to $S$}, i.e.
there exist an integer $m$ and a quasi-finite morphism from $Y$ to a suitable
standard projective scheme $\mathbf P_S^N$ such that
$(\mathcal L')^{\otimes m}$ is isomorphic to the inverse image of
$\mathcal O_{\mathbf P_S^N}(1)$.

Notice that even if $X$ is separated over $S$, it may happen that $Y$ is not
separated over $S$ (a situation occurring in entirely nonpathological moduli
problems). If $1^\circ$ holds, then $Y$ is separated if and only if the
equivalence relation defined by $G$ has a closed graph, i.e. if
$G\times X\to X\times X$ has closed image (it is then a closed immersion). If
$Y$ is separated, then $\mathcal L'$ is pre-ample on $Y$ relative to $S$ if
and only if it is ample, i.e. if a suitable tensor power defines a projective
immersion. In the moduli problems mentioned in the introduction, one can show
that the resulting equivalence relation does indeed have a closed graph.

\medskip
\noindent\textbf{REMARKS 8.1. ---} We have assumed
$G=\operatorname{GP}(n)_S$ to fix ideas and because this is presently the most
important case in practice. A reasonable hypothesis on $G$ would instead seem
to be that $G$ is one of the ``forms'' over $S$ of one of the Tohokû groups
(whose construction over the integers was also carried out by CHEVALLEY). The
only positive fact known to me in the direction of the preceding conjecture
is the following: let $X$ be \emph{an affine scheme over a field $k$ of
characteristic $0$, on which the group $\operatorname{Gl}(n)_k$ or
$\operatorname{GP}(n-1)_k$ acts freely. Then the equivalence relation defined
by $G$ is effective, the quotient $X/G$ is also affine, and the morphism
$X\to X/G$ is flat and surjective}. The proof uses the following fact (which
at present has been proved only in characteristic $0$): when $G$ is made to
act on the affine coordinate ring $A$ of $G$, regarded as a vector space over
$k$, the trivial representation of $G$ occurs only once (in a composition
series of a finite-dimensional $G$-stable vector subspace over $k$). It seems
possible that a systematic use of
\src{118}
the theory of linear representations of $G$ will eventually provide a proof
of the conjecture, at least over a base field. When the base is no longer a
field, the speaker knows nothing.

\section*{BIBLIOGRAPHY}
\addcontentsline{toc}{section}{Bibliography}

\begin{description}
\item[{[1]}] GROTHENDIECK (Alexander). --- Techniques de descente et théorèmes
d'existence en géométrie algébrique, I: Généralités, Séminaire Bourbaki,
1959/60, no. 190, 29 pp.

\item[{[2]}] GROTHENDIECK (Alexander). --- Techniques de descente et théorèmes
d'existence en géométrie algébrique, II: Le théorème d'existence en théorie
formelle des modules, Séminaire Bourbaki, 1959/60, no. 195, 22 pp.

\item[{[3]}] GROTHENDIECK (A.) and DIEUDONNÉ (J.). --- Éléments de géométrie
algébrique, I: Le langage des schémas. --- Paris, Presses universitaires de
France, 1960 (Institut des hautes Études scientifiques, Publications
mathématiques, 4).
\end{description}

\par\smallskip
\begingroup\footnotesize
\endgroup
